\documentclass[fleqn,12pt,a4paper,twoside]{article}

\usepackage{graphicx,calc,soul,amsmath,fancyhdr,psfrag,tikz}
\usepackage{enumitem}
\usepackage{listings}
\usepackage{float}
\usetikzlibrary{shapes.geometric,arrows,decorations.markings}

% Set the dimensions of the paper.
\setlength{\paperwidth}{210mm}
\setlength{\paperheight}{297mm}

% Set the dimensions of the text.
\setlength{\textwidth}{\paperwidth-6cm}
\setlength{\textheight}{\paperheight-2in}

\setlength{\headheight}{15pt}
\setlength{\marginparwidth}{3cm - \marginparsep -0.25in}

% Adjust margins to fit the paper and text dimensions.
\setlength{\oddsidemargin}{0.5\paperwidth-0.5\textwidth-1in}
\setlength{\evensidemargin}{\oddsidemargin}
\setlength{\topmargin}{0.5\paperheight-0.5\textheight-1in-\headheight-\headsep}

\markright{Experiment 3F1: Flight Control}

\begin{document}
\thispagestyle{empty}

\begin{center}
        \bfseries
        \rule{\textwidth}{.8mm} \\[.3cm]
        ENGINEERING TRIPOS PART IIA\\[1.5cm]
        ELECTRICAL AND INFORMATION ENGINEERING\\
        TEACHING LABORATORY\\[1.5cm]
        MODULE EXPERIMENT 3F1 \\[1.5cm]
        FLIGHT CONTROL\\[1.5cm]

        Name: Liu Zihe \\
        CRSID: zl559 \\
        College: Homerton \\
        Date: 29 November 2025 \\[1.5cm]
        \rule{\textwidth}{.8mm}\\[.5cm]
\end{center}
\vfill
\subsection*{Objectives:}

\begin{itemize}
        \item Simulation of various aircraft models on the computer.
        \item Study real-time (manual) control and the limitations imposed
              by time delays.
        \item Design of a simple autopilot.
        \item Illustrate frequency response concepts in analogue and digital control
              systems, conditions for oscillation in feedback systems and stability.
        \item Gain familiarity with Python and Jupyter Notebooks.
\end{itemize}
\vfill\vfill\null

\newpage

\section*{3F1 Lab Worksheet}

\begin{tabbing}
        \textbf{2.1 Simplified Aircraft Model} \\[3mm]
        Transfer function $= \dfrac{N}{s^2 + Ms} = \dfrac{10}{s^2 + 10s}$ \\[3mm]
        \verb+num = [10]+ \hspace*{2cm} \verb+den = [1, 10, 0]+\\
\end{tabbing}

\begin{tabbing}
        \textbf{2.2 Modelling Manual Control}\\[3mm]
        Controller transfer function $= K(s) = ke^{-Ds}$\\[3mm]
        $k= 0.8$ \hspace*{2cm} $D= 1$ s\\[3mm]
        Phase margin $= 45$ degrees\\[3mm]
        Amount of extra time delay which can be tolerated $= \dfrac{45 \times \pi/180}{0.8} \approx 1.0$ s\\
\end{tabbing}

\begin{tabbing}
        \textbf{2.3 Pilot Induced Oscillation}\\[3mm]
        Period of oscillation (observed) $= 5$ s\\[3mm]
        Period of oscillation (theoretical) $= \dfrac{2\pi}{\omega_{pc}} = \dfrac{2\pi}{0.8} \approx 7.85$ s\\
\end{tabbing}

\begin{tabbing}
        \textbf{2.4 Sinusoidal disturbances}\\[3mm]
        \= Maximum stabilising gain $= 4.5$\\[3mm]
        \> Gain at 0.66 Hz $= 5.5$ dB \hspace*{2cm} \= Phase at 0.66 Hz $= -90$ degrees\\[1mm]
\end{tabbing}

\begin{tabbing}
        \textbf{2.5 Unstable Aircraft}\\[3mm]
        Fastest pole at $T= 0.6$\\[3mm]
\end{tabbing}

\begin{tabbing}
        \textbf{3.2 Autopilot with Proportional Control}\\[3mm]
        \= Proportional gain $K_c= 12$ \hspace*{2cm} Period of oscillation $T_c= 2$ s\\
\end{tabbing}

\begin{tabbing}
        \textbf{3.3 Autopilot with PID Control}\\[3mm]
        Transfer function of PID controller $= K_p \left(1 + \dfrac{1}{T_i s} + T_d s\right)$\\[3mm]
        PID constants: $K_p = 7.2 \hspace*{1cm} T_i = 1.0 \hspace*{1cm} T_d = 0.25$\\[4mm]
        Adjusted value of \,$T_d = 0.35$\\
\end{tabbing}

\begin{tabbing}
        \textbf{3.4 Integrator Wind-up}\\[3mm]
        Integrator bound $Q= \dfrac{d \cdot T_i}{K_p} = \dfrac{2 \times 1.0}{7.2} = \dfrac{5}{18} \approx 0.278$\\
\end{tabbing}

\pagebreak

\section*{3F1 Flight Control Lab Report}

\vspace{1cm}

\noindent \textbf{1.} (§2.2 Modelling Manual Control) Nyquist diagram (from Bode diagram) for controller in series with plant.

The open-loop transfer function is $K(s)G(s) = 0.8 e^{-s} \cdot \dfrac{10}{s^2 + 10s}$.

From the Bode diagram (Figure 2):
\begin{itemize}
        \item At low frequencies ($\omega \to 0$): magnitude $\to \infty$, phase $\to -90°$
        \item At gain crossover ($\omega_{gc} = 0.8$ rad/s): magnitude $= 0$ dB, phase $\approx -135°$
        \item Phase crosses $-180°$ at approximately $\omega = 1.5$ rad/s
\end{itemize}

\begin{center}
        \begin{tikzpicture}[scale=1.1]
                \draw[->] (-3,0) -- (3,0) node[right] {Re};
                \draw[->] (0,-3) -- (0,3) node[above] {Im};
                % Draw Nyquist curve - starts from negative Im axis, spirals inward
                \draw[thick,blue,
                        decoration={markings, mark=at position 0.3 with {\arrow{>}}, mark=at position 0.7 with {\arrow{>}}},
                        postaction={decorate}]
                plot[smooth] coordinates {(1.5,-2.8) (1.2,-2) (0.8,-1.2) (0.3,-0.7) (0,-0.5) (-0.3,-0.3) (-0.5,-0.1) (-0.4,0.1) (-0.3,0.3) (0,0.5) (0.3,0.7) (0.8,1.2) (1.2,2) (1.5,2.8)};
                % Mark key points
                \node[font=\footnotesize] at (2,-2.5) {$\omega \to 0$};
                \node[font=\footnotesize] at (-0.8,0) {$\omega_{gc}$};
                \node[font=\footnotesize] at (2,2.5) {$\omega \to 0^-$};
                \filldraw[red] (-1,0) circle (2pt) node[below,font=\footnotesize] {$-1$};
                \node[align=center,font=\footnotesize] at (0,-3.5) {Nyquist diagram of $K(s)G(s)$};
        \end{tikzpicture}
\end{center}

\noindent \textbf{2.} (§2.2 Modelling Manual Control) Are you using any integral action? Give a brief explanation. What does this imply about the accuracy of the model of the human controller?

\vspace{8cm}

\pagebreak

\noindent \textbf{3.} (§2.3 Pilot Induced Oscillation) Explain the oscillation of the feedback loop.
How does your observed period of oscillation compare to the theoretical prediction?

\vspace{5cm}

\noindent \textbf{4.} (§2.3 Pilot Induced Oscillation) Can you give a rough guideline to the control designer to make PIO less likely?

\vspace{5cm}

\noindent \textbf{5.} (§2.4 Sinusoidal Disturbances) Was your manual input able to reduce the error (as compared to providing no input)?

\vspace{5cm}
\pagebreak

\noindent \textbf{6.} (§2.5 An Unstable Aircraft) Nyquist diagram for $G_2(s).$
\begin{center}
        \begin{tikzpicture}[scale=1.3]
                \draw[->] (-2,0) -- (2,0) node[right] {Re};
                \draw[->] (0,-2) -- (0,2) node[above] {Im};
                \node[align=center,font=\footnotesize] at (0,-2.5) {Nyquist diagram of $G_2(s)$};
        \end{tikzpicture}
\end{center}

\noindent \textbf{7.} (§2.5 An Unstable Aircraft) Explain, using the Nyquist criterion, why the feedback system is stable with a proportional gain greater than 0.5.

\vspace{5cm}

\noindent \textbf{8.} (§2.5 An Unstable Aircraft) Sketch of a Nyquist diagram for $G_2(s).$ with a small time delay $D$.
\begin{center}
        \begin{tikzpicture}[scale=1.3]
                \draw[->] (-2,0) -- (2,0) node[right] {Re};
                \draw[->] (0,-2) -- (0,2) node[above] {Im};
                \node[align=center,font=\footnotesize] at (0,-2.5) {Nyquist diagram with small time delay};
        \end{tikzpicture}
\end{center}

\pagebreak

\noindent \textbf{9.} (§3.4 Integrator Wind-up) Explain how you calculated the bound on $Q$.


\end{document}
